% glosario academico
\thispagestyle{empty}
\PageDecor{5}
\SectionTag{SECCIÓN INFORMATIVA}
\DocHeading{Glosario Académico}
\DocQuestion{¿Qué palabras necesito comprender para desenvolverme en una institución de nivel superior?}
\DocItem{Hora cátedra: uno o medio módulo de clases (35 o 40 min.) - Espacio curricular: materia-asignatura - Cátedra compartida o equipo de cátedra.}
\DocItem{Parcial - Final - Promoción directa - Coloquio - Regularidad: se refiere a la condición de alumno en relación a la asistencia a clases - Homologación - Correlatividades - Currículum - Práctica - Talleres - Seminarios - Exámenes finales: instancias de evaluación final sobre un espacio curricular (febrero-marzo; julio-agosto; noviembre-diciembre) - 1er, 2do, 3er llamado: instancias de mesas de exámenes finales - Programa - Proyecto: consiste en la planificación del espacio curricular que el docente realiza para el ciclo lectivo (objetivos, contenidos, estrategias didácticas, recursos, evaluación).}
\DocPara{La vida académica está regulada por un conjunto de normativas (leyes, decretos, resoluciones, disposiciones, circulares) con alcance a nivel nacional, provincial e institucional.}
\DocQuestion{¿Qué es una Normativa?}
\DocPara{Una normativa es un documento que contiene reglas, obligaciones, derechos y garantías. En este caso vienen a ordenar la vida en las instituciones y a regular la actividad que desarrollan los sujetos que las habitan, quienes suponen ciertos consensos para vivir en comunidad. Pueden ser disposiciones, resoluciones, decretos, leyes, estatutos, acuerdos, entre otros.}
